\documentclass[xcolor=monochrome,11pt]{beamer}

% ---------- Thème simple, noir et blanc ----------
\usetheme{default}
\usecolortheme{default}
\setbeamertemplate{navigation symbols}{}
\setbeamertemplate{footline}[frame number]
\setbeamercolor{block title}{fg=black,bg=white}
\setbeamercolor{block body}{fg=black,bg=white}

\usepackage[utf8]{inputenc}
\usepackage[T1]{fontenc}
\usepackage[french]{babel}
\usepackage{amsmath}

\title{La création d'un livre au prisme des écrans et de leurs mondes~:\\
une perspective post-phénoménologique}
\subtitle{Projet de recherche --- Session Été 2026}
\author{Guizaoui Yann \\ \small Matricule 20321059}
\institute{Remis à Michael Sinatra \\ Université de Montréal \\ Département de littératures et de langues du monde}
\date{Date de remise : 17 août 2026}

\begin{document}

% ---------- Slide 1 : Titre ----------
\begin{frame}
  \titlepage
\end{frame}

% ---------- Slide 2 : Problématique ----------
\begin{frame}{Problématique}
  \begin{block}{Question de recherche centrale}
    De quelle manière la techno-esthéticité des dispositifs numériques
    transforme-t-elle l'expérience du livre ainsi que les formes que
    celui-ci est susceptible d'adopter~?
  \end{block}
\end{frame}

\begin{frame}{Genèse de la définition du livre}
  \footnotesize
  \begin{itemize}
    \item Point de départ~: mémoire sur la philosophie du jeu vidéo
      $\rightarrow$ intérêt pour l'impact du numérique sur les activités
      culturelles, en particulier la lecture.
    \item Constat~: la littérature existante porte surtout sur les livres
      numériques \emph{homothétiques} (fidèles au codex) et néglige les
      fictions interactives textuelles.
    \item Déplacement de la question~: non plus « que fait le numérique au
      livre ? » mais « qu'est-ce qu'un livre numérique ? ».
    \item Recours à la phénoménologie de Poulet~: le livre comme artefact
      qui abolit la frontière dedans/dehors, lecteur/texte.
    \item Croisement avec les théories simulationnistes de la mémoire
      (Michaelian)~: se souvenir, c'est simuler le passé.
  \end{itemize}
\end{frame}

\begin{frame}{Une définition du livre}
  \footnotesize
  Le livre~: une certaine structure textuelle mnémonique dont le contenu
  s'actualise à la fois par la médiation d'un simulateur externe (la
  cognition du lecteur) et par sa propre structure (codex, parchemin, vase,
  programme manifesté sur un moniteur), ayant la capacité d'abolir pour le
  lecteur la frontière entre son propre espace-temps et celui qui est
  simulé.
  \vspace{0.3cm}
  \begin{itemize}
    \item Condition~: le texte doit demeurer le vecteur principal de la
      simulation --- il suffit d'observer, lire et interagir avec lui.
    \item Conséquence~: le livre n'est plus un artefact prénumérique figé ;
      le numérique peut produire ses \emph{propres} livres --- ce sont eux
      qui mériteraient l'appellation de « livre numérique ».
  \end{itemize}
\end{frame}

% ---------- Slide 3 : Cadre théorique ----------
\begin{frame}{Justification de la techno-esthétique et du cadre théorique}
  \footnotesize
  \begin{itemize}
    \item Double formation~: philosophie (esthétique) et sciences
      cognitives (interaction humain-machine) $\rightarrow$ cadre
      théorique croisant \textbf{post-phénoménologie} et
      \textbf{théorie des 4E} (embedded, extended, embodied, enactive).
    \item Lien historique~: les fondateurs de la théorie des 4E sont
      souvent issus de la philosophie.
    \item Post-phénoménologie~: étudier un phénomène technologique sans
      a priori $\rightarrow$ études empiriques $\rightarrow$ analyse
      conceptuelle (épistémologique, politique, esthétique, métaphysique).
    \item Les sciences cognitives fournissent le réservoir d'études
      empiriques (co-agentivité, perception du corps en espace virtuel...)
      et une passerelle vers le design (UX research).
  \end{itemize}
\end{frame}

\begin{frame}{Pourquoi « techno-esthétique » ?}
  \footnotesize
  \begin{itemize}
    \item Référence à Simondon~: intérêt pour ce que le dispositif
      technique a à offrir, pas seulement pour ce qu'on peut en tirer
      philosophiquement.
    \item Notion d'\emph{individu technique}~: une chose qui se conserve
      dans son devenir en augmentant sa stabilité à chaque changement
      $\rightarrow$ critère pour discriminer ce qui peut être généralisé.
    \item Simondon envisageait un projet de « techno-esthétique » /
      « esthético-technique » jamais développé --- repris ici comme
      \emph{adjectif}~: ce qui compte, c'est le dispositif \emph{mis en
      relation} (usage, contemplation...).
    \item Cette expérience est plurielle $\rightarrow$ modes
      d'expériences techno-esthétiques du numérique.
    \item Esthétique~: accès à la connaissance par des sens eux-mêmes
      affectés par le milieu $\rightarrow$ intérêt d'une lecture esthétique
      plutôt qu'ontologique des dispositifs numériques.
  \end{itemize}
\end{frame}

% ---------- Slide 4 : Annonce des démarches ----------
\begin{frame}{Annonce des démarches}
  Trois démarches structurent le projet, chacune correspondant à un
  chapitre~:
  \begin{enumerate}
    \item Des dispositifs numériques à leurs modes d'expression
    \item Cartographie des fictions textuelles interactives
    \item Le livre-monde~: spatialisation, simulation et autopoïèse
      narrative
  \end{enumerate}
\end{frame}

% ---------- Slide 5 : Démarche 1 ----------
\begin{frame}{Démarche 1 --- Des dispositifs numériques à leurs modes d'expression}
  \begin{block}{Problématique}
    \small Quels modes d'expression peut-on dégager de la structure
    logique des dispositifs numériques~?
  \end{block}
\end{frame}

\begin{frame}{1. L'architecture logique du microprocesseur}
  \footnotesize
  \begin{itemize}
    \item Dispositif technique~: architecture de Von Neumann envisagée
      comme système en \emph{individuation} (métastable, non figé).
    \item « Avant d'arriver à la techno-esthétique d'un ensemble, il faut
      considérer celle de l'individu » (Simondon).
    \item Analyse à la fois formelle (fonctionnement logique) et
      matérielle (organisation physique des opérations).
  \end{itemize}
\end{frame}

\begin{frame}{2. La boucle de rétroaction utilisateur--machine}
  \footnotesize
  \begin{itemize}
    \item Dispositif psycho-technique~: un ordinateur fonctionne en boucle
      ouverte (entrées / sorties), via des interfaces (clavier,
      moniteur...).
    \item Étude de cas~: évolution du moniteur entre 1951 et 1977.
    \item Comment cette évolution a transformé le régime d'interaction
      humain-machine, jusqu'à intégrer l'état de la machine au sein du
      \emph{sensorium commune} de l'utilisateur.
  \end{itemize}
\end{frame}

\begin{frame}{3. Modes d'expression techno-esthétiques du numérique}
  \footnotesize
  Cinq types de propriétés, constituant la grille de lecture du corpus
  (chapitre~II)~:
  \begin{itemize}
    \item \textbf{Relationnelles}~: interactivité, ludogénéité
      (Vial, 2016), projectivité
    \item \textbf{Ontologiques}~: reproductibilité, procéduralité
    \item \textbf{Référentielles}~: para-authenticité (Lee, 2004),
      virtualité
    \item \textbf{Perceptives}~: spatialité, temporalité, aoristicité
      (espace fini sans bords, perçu comme infini)
  \end{itemize}
\end{frame}

% ---------- Slide 6 : Démarche 2 ----------
\begin{frame}{Démarche 2 --- Cartographie des fictions textuelles interactives}
  \begin{block}{Problématique}
    \small Jusqu'où les fictions textuelles interactives ont-elles
    exploré les modes d'expression identifiés~?
  \end{block}
\end{frame}

\begin{frame}{Les fictions interactives textuelles (1930--2026)}
  \footnotesize
  \begin{itemize}
    \item L'interface reste celle du papier jusqu'en 1977~:
      \emph{Un livre dont vous êtes le héros} comme remédiation des
      logiques computationnelles vers le format livresque hérité de
      \emph{Donjons et Dragons}.
    \item Panorama des logiciels et langages majeurs~: Twine, Adrift,
      Ink, etc.
    \item Distinction fictions à choix / fictions procédurales, illustrée
      par le cas de \emph{SillyTavern}.
  \end{itemize}
\end{frame}

\begin{frame}{Typologie et \emph{terra incognita} du livre numérique}
  \footnotesize
  \begin{itemize}
    \item Production d'une typologie des fictions interactives textuelles
      selon leurs modalités de design, leurs types d'interaction et leur
      rapport au temps.
    \item Croisement avec les propriétés techno-esthétiques du chapitre~I
      pour localiser les espaces de design inexplorés.
    \item Deux hypothèses préliminaires~:
      \begin{itemize}
        \item une tridimensionnalité du texte absente du corpus existant ;
        \item des enjeux de design posés par les fictions interactives
          textuelles génératives par IAg.
      \end{itemize}
  \end{itemize}
\end{frame}

% ---------- Slide 7 : Démarche 3 ----------
\begin{frame}{Démarche 3 --- Le livre-monde}
  \begin{block}{Problématique}
    \small Dans quelle mesure la spatialisation et le récit
    quasi-autopoïétique constituent-ils un livre-monde, et à quelles
    conditions ce fantasme est-il réalisable~?
  \end{block}
\end{frame}

\begin{frame}{Le livre comme simulation}
  \footnotesize
  \begin{itemize}
    \item Articulation de la conception phénoménologique du livre chez
      Poulet avec la narratologie cognitive (\emph{storyworld}, Herman)~:
      le lecteur reconstruit un monde à partir des indices du texte.
    \item Mise en tension avec le débat continuiste / discontinuiste
      mémoire--imagination (Michaelian).
    \item Thèse défendue~: le livre numérique peut être pensé comme un
      simulateur, et non seulement comme le déclencheur de notre
      simulateur mental --- coupler un simulateur textuel à notre propre
      faculté de simulation du monde.
  \end{itemize}
\end{frame}

\begin{frame}{Spatialisation et générativité}
  \footnotesize
  \begin{itemize}
    \item Hypothèse~: la spatialisation 3D et la génération de texte
      restent peu exploitées dans les fictions interactives.
    \item Cause probable~: l'exigence de continuité textuelle --- en
      espace 2D/3D peu codifié, le texte s'organise en réseau (carte
      mentale), rendant la rédaction difficilement soutenable.
    \item La génération de texte rend cette continuité de nouveau
      envisageable.
    \item Réflexion sur un livre hypothétique~: un \emph{canvas} infini
      tridimensionnel où le lecteur se déplace librement, les connexions
      textuelles entre points d'intérêt étant générées à la volée.
  \end{itemize}
\end{frame}

\begin{frame}{Un éditeur de livre-monde}
  \footnotesize
  \begin{itemize}
    \item Présentation, sous forme de \emph{wireframes}, du prototype
      d'un éditeur permettant la conception de livres-mondes.
    \item Modalités de développement et principales fonctionnalités.
    \item Exemples de design justifiés (UX design), non fonctionnels,
      illustrant les choix envisagés.
  \end{itemize}
\end{frame}

% ---------- Slide 8 : Conclusion / Objectifs ----------
\begin{frame}{Conclusion / Objectifs}
  \begin{itemize}
    \item Une contribution \textbf{critique} sur la doxa du codex
    \item Une contribution \textbf{conceptuelle} sur la
      techno-esthéticité du numérique en général et du livre numérique en
      particulier
    \item Une contribution \textbf{typologique} identifiant les modes
      d'expression du numérique les plus mobilisés dans ces livres, ainsi
      que ceux qui demeurent marginaux ou absents
    \item Une contribution \textbf{prospective} esquissant les conditions
      de possibilité de formes de livres numériques non encore envisagées
  \end{itemize}
\end{frame}

% ---------- Slide 9 : Étude de cas technique ----------
\begin{frame}{Étude de cas technique~: API HTML-in-Canvas}
  \footnotesize
  \begin{itemize}
    \item Juin 2025~: implémentation de HTML-in-Canvas dans Three.js, en
      parallèle d'un projet personnel de livre numérique en espace 3D
      sans bord.
    \item Idée initiale~: afficher le texte à l'écran et actualiser son
      CSS selon la position de la caméra.
      \begin{itemize}
        \footnotesize
        \item Fonctionnait visuellement, mais trahissait le projet~:
          illusion de 3D, sans réel positionnement spatial du texte.
      \end{itemize}
    \item Problème en VR/WebXR~: seul le contexte WebGL apparaît en
      session WebXR $\rightarrow$ le texte (DOM~2D superposé) devient
      invisible.
    \item Solution explorée~: conversion du rendu HTML en texture SVG
      appliquée à un mesh 3D --- fonctionne, mais perte de
      l'accessibilité native du HTML/CSS.
    \item L'API HTML-in-Canvas répond directement à ce problème.
    \item Prochaine étape~: reprise du projet avec cette API et le
      polyfill de repalash.
  \end{itemize}
\end{frame}

\end{document}
